\chapter{前言}

\section*{为什么写这本书}

算法组新队员入门定位与导航，最常见的困境不是“找不到资料”，而是“资料太多且互相矛盾”：有人推荐先看《视觉 SLAM 十四讲》，有人让你直接读 FAST-LIO 源码；有人说定位就是卡尔曼滤波，有人说是因子图优化；有人把 Cartographer 当标准答案，有人只用 LIO。结果往往是：视频看了几十个小时，代码 clone 了一堆，真到机器人上还是不知道从哪下手。

这本书试图解决这个问题。它不假设你学过机器人学，也不假设你会 ROS，而是从“机器人怎么知道自己在哪”这个最朴素的问题出发，一层一层把概念、公式、算法、工程串起来。我们尽量做到三件事：
\begin{enumerate}
  \item \textbf{讲清“为什么”}：每个算法先说出它要解决什么问题、在前一个方法上改进了什么，再给公式；
  \item \textbf{推导到能算数}：关键公式从定义开始推，不跳步，并且尽量给一个可以手算的数字例子；
  \item \textbf{落到工程}：每个模块都说明输入输出、频率、参数、失效场景，而不是停在论文层面。
\end{enumerate}

\section*{本书面向谁}

主要读者是 RoboMaster 算法组的新队员，以及任何“会 C++、学过高数线代概率、但没系统学过定位导航”的同学。我们假设你：
\begin{itemize}
  \item 能看懂基本的矩阵运算、求导、概率分布；
  \item 写过一些 C++/Python，能跑通 Linux 命令；
  \item 对机器人有热情，愿意动手。
\end{itemize}
如果你已经做过一两个赛季的定位，这本书的后半部分（SLAM 篇、导航篇、工程篇）仍然值得一读——它会把你在调参中积累的“直觉”补上数学解释。

\section*{怎么读}

全书五篇三十章，建议按顺序读完第一篇（第 1--3 章），之后按需选择：

\begin{itemize}
  \item \textbf{一周速成}：第 1、2、4、6、11、16、17 章。建立“里程计—前端—后端—LIO”的主干。
  \item \textbf{建图方向}：第 1--3、7、11--15 章。重点是 Cartographer、地图表示与回环。
  \item \textbf{导航方向}：第 1、10、11、18--22 章。重点是 Nav2 与 CMU 3D 导航栈。
  \item \textbf{落地实战}：第 23、25--29 章。Z-Localization 的架构、ROS 2 工程、标定与调参。
\end{itemize}

每章开头有“本章导读”，结尾有“本章小结”“思考与练习”“延伸阅读”。建议至少把小结和练习过一遍——练习里的题目大多是“把本章公式代入具体数字”，能算出来才算真懂。

\section*{关于“前端/后端”}

这是本书反复出现的两个词，也是新手最容易混淆的地方。第 1 章先给直觉，第 11 章给出正式定义，并特别用一节讲清楚它们在三种语境下的不同含义：
\begin{itemize}
  \item SLAM 语境：前端造约束，后端解约束；
  \item 定位系统语境：前端是里程计，后端是先验地图配准与因子图；
  \item 导航系统语境：前端是感知（定位+建图），后端是决策（规划+控制）。
\end{itemize}
读任何一章时，请先确认该章用的是哪种语境。第 23 章 Z-Localization 用的是第二种，第 18 章 Nav2 用的是第三种。

\section*{数学约定}

为了让全书符号一致，我们约定：
\begin{itemize}
  \item 向量、矩阵用粗体：$\bm{p},\bm{R}$；标量用普通斜体：$x,\theta$；
  \item 估计值加帽子：$\hat{\bm{x}}$；IMU 原始测量加波浪号：$\tilde{\bm{\omega}}$；
  \item 坐标系：$W$ 世界系、$O$ 里程计系、$B$ 机体系、$L$ 雷达系、$C$ 相机系；
  \item 位姿 $T_{AB}$ 表示“从 $B$ 系到 $A$ 系”的变换，满足 ${}^{A}\bm{p}=T_{AB}\,{}^{B}\bm{p}$；
  \item 重力 $\bm{g}=[0,0,-9.81]^{\top}$，世界系 $z$ 轴向上；
  \item 旋转用右扰动更新：$\bm{R}\leftarrow\bm{R}\exp(\delta\bm{\phi}^{\wedge})$。
\end{itemize}
完整的符号表见第 2 章与附录 A。

\section*{关于开源代码}

本书大量参考了以下开源项目，特此致谢：
\begin{itemize}
  \item \textbf{FAST-LIO / FAST-LIO2 / Point-LIO / VoxelMap}（香港大学 MARS 实验室，Fu Zhang 组）：第 16--17 章的主要素材；
  \item \textbf{Cartographer}（Google）：第 15 章；
  \item \textbf{Nav2}（Open Navigation LLC 与社区）：第 18、22 章；
  \item \textbf{CMU Autonomous Exploration Development Environment}（CMU AirLab，Ji Zhang 组）：第 21 章；
  \item \textbf{Z-Localization}（东北大学 T-DT 战队，作者 mozihe）：第 23 章。该项目的设计说明写得非常坦诚，本章的很多工程判断直接来自作者的分享。
\end{itemize}
书中对代码的描述基于公开文档与论文，若与最新版本有出入，请以官方仓库为准。

\section*{致谢与反馈}

感谢算法组历届队员积累的调试笔记与踩坑记录，它们是本书“工程篇”的素材来源。感谢所有把代码与设计思路开源出来的团队——没有这些开源项目，这本书无从写起。

本书难免有错误与过时之处（尤其涉及具体参数与版本的部分）。如果你发现公式推导有误、代码描述过时、或者某一段实在看不懂，欢迎指出。对一本教材来说，读者的困惑和作者的笔误同样重要。
